% القسم: المناهج
% الفصل: الاستقراء
% المبحث: الاستقراء على الأعداد الطبيعية

\documentclass[../../../include/open-logic-section]{subfiles}

\begin{document}

\olfileid{mth}{ind}{inN}

\olsection{الاستقراء على~$\Nat$}

الاستقراء، في أبسط صوره، تقنية لإثبات نتائج لجميع الأعداد الطبيعية.
وهو يستعمل حقيقة أننا، إذا بدأنا من $0$ وأضفنا~$1$ مراراً، نصل في
النهاية إلى كل عدد طبيعي. ولذلك، لكي نثبت صدق شيء لكل عدد، نستطيع:
(1)~أن نثبت صدقه للعدد $0$، و(2)~أن نبين أنه كلما صدق لعدد~$n$ صدق
أيضاً للعدد التالي~$n+1$. فإذا اختصرنا «للعدد~$n$ الخاصية $P$» إلى
$P(n)$، و«للعدد~$k$ الخاصية $P$» إلى $P(k)$، وهكذا، فإن البرهان
بالاستقراء على $P(n)$ لكل $n \in \Nat$ يتألف من:
\begin{enumerate}
\item برهان على $P(0)$، و
\item برهان على أنه، لكل~$k$، إذا كان $P(k)$ فإن $P(k+1)$.
\end{enumerate}
ولتوضيح ذلك تماماً، افترض أن لدينا كلاً من (1) و(2). تخبرنا (1) أن
$P(0)$ صادقة. وإذا كانت لدينا (2) أيضاً، نعلم بوجه خاص أنه إذا كانت
$P(0)$ صادقة فإن $P(0+1)$، أي $P(1)$، صادقة. وهذا ينتج من القضية
العامة «لكل~$k$، إذا كان $P(k)$ فإن $P(k+1)$» بإحلال~$0$ محل~$k$.
فنحصل، بالقياس الشرطي، على $P(1)$. ومن (2) مرة أخرى، مع إحلال $1$ هذه
المرة محل~$k$، نحصل على: إذا كان $P(1)$ فإن~$P(2)$. ولما كنا قد
أثبتنا~$P(1)$ للتو، نحصل بالقياس الشرطي على~$P(2)$. وهكذا. فلكل
عدد~$n$ نصل في النهاية، بعد فعل ذلك $n$ مرة، إلى~$P(n)$. ومن ثم تثبت
(1) و(2) معاً $P(n)$ لأي $n \in \Nat$.

لننظر في مثال. افترض أننا نريد معرفة عدد المجاميع المختلفة التي يمكن
رميها بـ$n$ من أحجار النرد. ومع أن الأمر قد يبدو سخيفاً، فلنبدأ بـ$0$
من أحجار النرد. إذا لم يكن لديك أي حجر نرد، فليس ثمة إلا مجموع واحد
ممكن أن «ترميه»: لا نقاط على الإطلاق، ومجموعها~$0$. إذن عدد الرميات
الممكنة المختلفة هو~$1$. وإذا كان لديك حجر واحد فقط، أي $n=1$، فهناك
ست قيم ممكنة من $1$ إلى~$6$. وبحجرين يمكن رمي أي مجموع من $2$ إلى
$12$، أي $11$ احتمالاً. وبثلاثة أحجار يمكن رمي أي عدد من $3$ إلى
$18$، أي $16$ احتمالاً مختلفاً. تبدو الأعداد $1$ و$6$ و$11$ و$16$
نمطاً: فهل الجواب $5n+1$؟ وبالطبع $5n+1$ هو الحد الأقصى الممكن، إذ لا
يوجد إلا $5n+1$ عدداً بين $n$، وهو أقل قيمة يمكن رميها بـ$n$ من أحجار
النرد عندما تكون كلها $1$، و$6n$، وهو أعلى قيمة عندما تكون كلها $6$.

\begin{thm}
  يمكن بـ$n$ من أحجار النرد رمي جميع القيم الممكنة وعددها $5n+1$ بين
  $n$ و$6n$.
\end{thm}

\begin{proof}
لتكن $P(n)$ القضية: «يمكن رمي أي عدد بين $n$ و~$6n$ باستعمال $n$ من
أحجار النرد». ولاستعمال الاستقراء نثبت:
\begin{enumerate}
\item \emph{أساس الاستقراء} $P(0)$، أي إنه من دون أحجار نرد لا يمكن
  رمي إلا المجموع $0$.
\item \emph{خطوة الاستقراء}: لكل $k$، إذا كان $P(k)$ فإن~$P(k+1)$.
\end{enumerate}

تثبت (1) لأن الرمية الخالية من أحجار النرد لها مجموع واحد هو $0$، وهو
القيمة الوحيدة بين $0$ و~$6\cdot0$.

ولإثبات (2)، نفترض مقدم القضية الشرطية، أي~$P(k)$. ويسمى هذا الافتراض
\emph{فرض الاستقراء}. ونستعمله لإثبات~$P(k+1)$. والجزء الصعب هو إيجاد
طريقة للتفكير في القيم الممكنة لرمي $k+1$ من أحجار النرد بدلالة القيم
الممكنة لرميات $k$ منها مع رمية الحجر الإضافي رقم $k+1$؛ لكن هذا هو ما
ينبغي فعله إذا أردنا استعمال فرض الاستقراء.

يقول فرض الاستقراء إننا نستطيع الحصول على أي عدد بين $k$ و~$6k$
باستعمال $k$ من أحجار النرد. فإذا رمينا~$1$ بالحجر رقم $(k+1)$، أضاف
ذلك $1$ إلى المجموع. ولذلك نستطيع رمي أي قيمة بين $k+1$ و$6k+1$ برمي
$k$ من أحجار النرد ثم رمي~$1$ بالحجر رقم $(k+1)$. فما الذي بقي؟ القيم
من $6k+2$ إلى $6k+6$. ويمكن الحصول عليها برمي $6$ بكل واحد من الأحجار
$k$، ثم عدداً بين $2$ و$6$ بالحجر رقم $(k+1)$. ويعني هذا كله أننا
نستطيع بـ$k+1$ من أحجار النرد رمي أي عدد بين $k+1$ و$6(k+1)$؛ أي إننا
أثبتنا~$P(k+1)$ باستعمال الفرض~$P(k)$، وهو فرض الاستقراء.
\end{proof}

كثيراً ما نستعمل الاستقراء حين نريد إثبات شيء عن متوالية من الأشياء،
كالأعداد أو المجموعات، تكون هي نفسها معرّفة «استقرائياً»، أي بتعريف
الشيء رقم $(n+1)$ بدلالة الشيء رقم $n$. فمثلاً، نستطيع تعريف
المجموع~$s_n$ للأعداد الطبيعية حتى~$n$ كما يأتي:
\begin{align*}
  s_0 & = 0\\
  s_{n+1} & = s_n + (n+1)
\end{align*}
يعطينا هذا التعريف:
\begin{align*}
  s_0 & = 0,\\
  s_1 & = s_0 + 1 && = 1,\\
  s_2 & = s_1 + 2 && = 1 + 2 = 3\\
  s_3 & = s_2 + 3 && = 1 + 2 + 3 = 6, \text{ وهكذا.}
\end{align*}
ونستطيع الآن أن نثبت بالاستقراء أن $s_n = n(n+1)/2$.

\begin{prop}
  $s_n = n(n+1)/2$.
\end{prop}

\begin{proof}
  علينا أن نثبت (1) أن $s_0 = 0\cdot(0 + 1)/2$، و(2) أنه إذا كان
  $s_k = k(k+1)/2$ فإن $s_{k+1} = (k+1)(k+2)/2$. أما (1) فبديهية.
  ولإثبات (2)، نفترض فرض الاستقراء: $s_k = k(k+1)/2$. وعلينا باستعماله
  أن نبين أن $s_{k+1} = (k+1)(k+2)/2$.

  فما قيمة $s_{k+1}$؟ بحسب التعريف،
  $s_{k+1} = s_k + (k+1)$. وبحسب فرض الاستقراء،
  $s_k = k(k+1)/2$. ونستطيع إحلال ذلك في المعادلة السابقة، ثم لا نحتاج
  إلا إلى شيء من حساب الكسور:
  \begin{align*}
    s_{k+1} & = \frac{k(k+1)}{2} + (k+1)\\
    & = \frac{k(k+1)}{2} + \frac{2(k+1)}{2}\\
    & = \frac{k(k+1) + 2(k+1)}{2}\\
    & = \frac{(k+2)(k+1)}{2}.
  \end{align*}
\end{proof}

والدرس المهم هنا هو أنك إذا كنت تثبت شيئاً عن متوالية $a_n$ معرّفة
استقرائياً، فالاستقراء هو الطريق الواضح. وحتى إن لم يكن كذلك، كما في
حالة القيم الممكنة لرميات النرد، يمكنك استعمال الاستقراء إذا استطعت أن
تربط حالة~$k+1$ بحالة~$k$ على نحو ما.

\end{document}
